{\footnotesize
\begin{threeparttable}
\begin{tabular}{lrrrrrrrrr}
\toprule
 & \multicolumn{3}{c}{Within 6 months (\%)} & \multicolumn{3}{c}{Within 12 months (\%)} & \multicolumn{3}{c}{$\Delta \ln w$} \\
\cmidrule(lr){2-4}\cmidrule(lr){5-7}\cmidrule(lr){8-10}
Relative decline & G1 & G3 & Diff. & G1 & G3 & Diff. & G1 & G3 & Diff. \\
\midrule
All & 30.1 & 29.4 & 0.7 & 50.6 & 50.9 & $-$0.3 & $-$0.069 & $-$0.119 & 0.050 \\
5 or more & 29.1 & 27.6 & 1.5 & 49.3 & 48.8 & 0.6 & $-$0.086 & $-$0.144 & 0.058 \\
10 or more & 29.0 & 27.4 & 1.7 & 48.7 & 47.5 & 1.2 & $-$0.093 & $-$0.154 & 0.062 \\
15 or more & 29.4 & 27.1 & 2.3 & 48.8 & 46.6 & 2.2 & $-$0.088 & $-$0.151 & 0.063 \\
20 or more & 29.0 & 26.2 & 2.8 & 48.2 & 44.6 & 3.6 & $-$0.098 & $-$0.158 & 0.060 \\
\bottomrule
\end{tabular}
\begin{tablenotes}[flushleft]\footnotesize
\item \textit{Note:} Same sample as panel B of \hyperref[tab:escada]{Table~\ref*{tab:escada}}, excluding workers without a positive origin wage. Group 3 means are recomputed with Group 1 weights in the 16 cells of tenure and initial wage (\hyperref[eq:dfl_reweight]{Equation~(\ref*{eq:dfl_reweight})}), within each band of relative decline, so that the two groups have the same composition along these dimensions. The first two blocks give the share reemployed within 6 and 12 months, and the third the change in log wages. ``Diff.'' is in percentage points in the first two blocks and in log points in the third. Source: RAIS, authors' calculations.
\end{tablenotes}
\end{threeparttable}
}
